\begin{proof}[Proof of \cref{obj:lem:scalar-implicit-function}]
\begin{enumerate}
\item Since \(D\) is open and contains \((\mathbf b_0,\chi_0)\), the smoothness hypothesis gives continuous differentiability at this point. Define the continuous linear map
\[
L_\chi:\mathbb R\to\mathbb R,
\qquad
L_\chi(t)=DF(\mathbf b_0,\chi_0)[(\mathbf 0,t)].
\]
The chain rule for the map \(\chi\mapsto(\mathbf b_0,\chi)\) gives
\[
\left.\frac{d}{d\chi}F(\mathbf b_0,\chi)\right|_{\chi=\chi_0}
=L_\chi(1)\ne0.
\]
By linearity,
\[
L_\chi(t)=tL_\chi(1)
\qquad(t\in\mathbb R).
\]
Define
\[
J_\chi:\mathbb R\to\mathbb R,
\qquad
J_\chi(t)=\frac{t}{L_\chi(1)}.
\]
This map is continuous and linear, and the preceding identity gives
\[
L_\chi(J_\chi(t))=t,
\qquad
J_\chi(L_\chi(t))=t
\qquad(t\in\mathbb R).
\]
Thus \(L_\chi\) is invertible. % lean: CausalSmith.Stat.LogoddsLowsmoothFrontier.scalar_implicit_function

\item The general \(C^1\) implicit-function theorem for Banach-space products, applied with scalar-coordinate derivative \(L_\chi\), supplies a function
\[
q:\mathbb R^d\to\mathbb R
\]
such that
\[
q(\mathbf b_0)=\chi_0,
\qquad
q\text{ is continuously differentiable at }\mathbf b_0.
\]
Its local graph identity, together with \(F(\mathbf b_0,\chi_0)=0\), gives
\[
F(\mathbf b,\chi)=0
\quad\Longleftrightarrow\quad
q(\mathbf b)=\chi
\]
throughout a neighborhood of \((\mathbf b_0,\chi_0)\). % lean: CausalSmith.Stat.LogoddsLowsmoothFrontier.scalar_implicit_function

\item Intersect this neighborhood with \(D\). The product topology supplies neighborhoods \(U_1\subseteq\mathbb R^d\) of \(\mathbf b_0\) and \(V_1\subseteq\mathbb R\) of \(\chi_0\) satisfying
\[
U_1\times V_1
\subseteq
D\cap
\bigl\{(\mathbf b,\chi):
F(\mathbf b,\chi)=0
\Longleftrightarrow q(\mathbf b)=\chi\bigr\}.
\]
Choose open neighborhoods \(U_2\subseteq\mathbb R^d\) and \(V\subseteq\mathbb R\) with
\[
\mathbf b_0\in U_2\subseteq U_1,
\qquad
\chi_0\in V\subseteq V_1.
\]
Continuous differentiability of \(q\) at \(\mathbf b_0\) also supplies a neighborhood \(S\subseteq\mathbb R^d\) such that
\[
\mathbf b_0\in\operatorname{int}(S),
\qquad
q\text{ is continuously differentiable on }S.
\]
% lean: CausalSmith.Stat.LogoddsLowsmoothFrontier.scalar_implicit_function

\item Continuity of \(q\) at \(\mathbf b_0\), the identity \(q(\mathbf b_0)=\chi_0\), and openness of \(V\) imply
\[
\mathbf b_0\in\operatorname{int}\bigl(q^{-1}(V)\bigr),
\qquad
q^{-1}(V)=\{\mathbf b\in\mathbb R^d:q(\mathbf b)\in V\}.
\]
Hence we may choose an open set \(U\subseteq\mathbb R^d\) satisfying
\[
\mathbf b_0\in U
\subseteq U_2\cap S\cap q^{-1}(V).
\]
The inclusions established above yield
\[
U\times V\subseteq U_1\times V_1\subseteq D,
\]
while restriction to \(U\subseteq S\) gives continuous differentiability of \(q\) on \(U\), and
\[
q(\mathbf b)\in V
\qquad(\mathbf b\in U).
\]
Finally, for every \(\mathbf b\in U\) and \(\chi\in V\), the graph identity on \(U_1\times V_1\) gives
\[
F(\mathbf b,\chi)=0
\quad\Longleftrightarrow\quad
q(\mathbf b)=\chi
\quad\Longleftrightarrow\quad
\chi=q(\mathbf b).
\]
These are the required properties. % lean: CausalSmith.Stat.LogoddsLowsmoothFrontier.scalar_implicit_function
\end{enumerate}
\end{proof}
